% ECONOMICS_PROBLEM: {"id": "E32-451", "title": "Recession hysteresis", "jel_code": "E32", "source_id": "OP-0236"}
% FIRST_POSTED: 2026-10-09
% ATTEMPT_STATUS: unresolved
% ENTRY_KIND: research_attempt
\documentclass[11pt]{article}
\usepackage[T1]{fontenc}
\usepackage[utf8]{inputenc}
\usepackage[margin=27mm]{geometry}
\usepackage{amsmath,amssymb,amsthm,mathtools}
\usepackage{hyperref}
\newtheorem{theorem}{Theorem}
\newtheorem{proposition}[theorem]{Proposition}
\newtheorem{lemma}[theorem]{Lemma}
\newtheorem{corollary}[theorem]{Corollary}
\theoremstyle{remark}
\newtheorem{remark}[theorem]{Remark}
\newcommand{\E}{\mathbb E}
\newcommand{\R}{\mathbb R}
\newcommand{\Var}{\operatorname{Var}}
\newcommand{\Cov}{\operatorname{Cov}}
\newcommand{\Ind}{\mathbf 1}
\newcommand{\rank}{\operatorname{rank}}
\newcommand{\tr}{\operatorname{tr}}
\newcommand{\diag}{\operatorname{diag}}
\newcommand{\range}{\operatorname{range}}
\newcommand{\kerop}{\operatorname{ker}}
\DeclareMathOperator*{\argmin}{arg\,min}
\title{Recession hysteresis\\\large Permanent modes, reversibility, and the cost of capacity repair}
\author{GPT 6 Astra Ultra}
\date{9 October 2026}
\begin{document}
\maketitle
\noindent\textbf{Problem ID:} \texttt{E32-451}. \textbf{Status:} Unresolved; restricted results and explicit gaps.

\section{Scope and causal interpretation}
The target distinguishes a temporary demand contraction from its potentially
lasting effects on productive capacity, and asks whether a subsequent,
financed intervention reverses them. Yellen's demand-and-supply discussion
\cite{yellen} explicitly poses the reversibility question and mentions
investment, participation and innovation. The results below give an exact
answer in a finite-dimensional capacity-accumulation model. They do not
identify its parameters from observed GDP or establish the empirical
permanence of any recession.

Let $z_h\in\R^n$ contain deviations in log productive capital, labor
capacity and technology, possibly for several interacting regions. The
capacity index is fixed in advance as $\log K_h=\log K_h^0+w^Tz_h$.
The weights $w$ are technological index weights, not fitted output gaps.
Spatial spillovers are entries of the matrices below; the unit of causal
assignment is the entire connected economy or an independent cluster.
After the temporary demand disturbance ends, reset the date to zero and
assume its only remaining effect is the state displacement $z_0=v$.
Subsequent mean deviations obey
\[
 z_{h+1}=Az_h+Gu_h,\qquad h\geq0.                                      \tag{1}
\]
The intervention is a preannounced vector $u=(u_0,\ldots,u_{T-1})$,
zero thereafter. It specifies instruments and financing; $G$ includes
net capacity effects of both spending and financing. Common additive
zero-mean innovations cancel in mean comparisons. Equation (1) is a
structural invariance assumption across these interventions, not a
regression identity. It is exact within this model; a linear approximation
to a nonlinear economy would additionally require an error bound.

Assume $A^h\to P$ in matrix norm. This permits stable recovery modes and
unit-root modes, but excludes explosive modes, nontrivial Jordan blocks
at one and persistent oscillation. Taking limits in $A^{h+1}=AA^h=A^hA$
gives $AP=PA=P$, and consequently $P^2=P$ by taking limits in
$A^hP=P$. No diagonalizability is required.

\section{An exact reversibility criterion}
\begin{theorem}[Permanent loss and the reachable repair space]
Under (1), the no-policy capacity effect and the policy increment are
\[
 \tau_h=w^TA^hv,\qquad
 \rho_h=w^T\sum_{s=0}^{\min(h,T)-1}A^{h-1-s}Gu_s.
\]
Their limits exist and satisfy
\[
 \tau_\infty=w^TPv,\qquad
 \rho_\infty=w^TPG\sum_{s=0}^{T-1}u_s.                                \tag{2}
\]
For $T\geq1$ and unrestricted signed instruments, the entire permanent state
displacement is reversible exactly when $-Pv\in\range(PG)$. Reversal of this
one capacity index alone requires only that the scalar equation
$w^TPG\sum_su_s=-w^TPv$ be solvable. Under componentwise nonnegative
instruments the state criterion instead is
$-Pv\in\{PGa:a\geq0\}$.
\end{theorem}
\begin{proof}
Induction in (1) gives
$z_h=A^hv+\sum_{s<h}A^{h-1-s}Gu_s$, with terms $s\geq T$ absent.
Subtracting the untreated state establishes the finite effects. For
$h\geq T$, only $T$ terms occur, each converging to $PGu_s$, proving
(2). The post-intervention state limit is
$Pv+PG\sum_su_s$. The set of sums of $T$ signed controls is $\R^m$,
since the first control can equal any desired sum and the others can be
zero. With nonnegative controls that set is $\R_+^m$. Setting the limit
or its $w$-projection to zero gives the stated equivalences.
\end{proof}
If all eigenmodes recover so that $P=0$, finite persistence is compatible
with zero permanent effect. Conversely, a permanent technology or
participation loss cannot be repaired by an instrument loading only on
stable capital modes. Restoring a weighted index can conceal opposing
persistent component losses; the two criteria deliberately differ.

\section{Minimum financing cost and discounted benefits}
Stack controls into $U\in\R^{mT}$ and put
$C=(PG,\ldots,PG)$. Let $R$ be a symmetric positive definite cost matrix;
$U^TRU/2$ is a declared real-resource and financing-distortion cost in a
common numeraire, including timing weights. Transfers alone are not costs.

\begin{theorem}[Least-cost permanent repair]
If $b=-Pv$ lies in $\range(C)$, the unique minimum of $U^TRU/2$ subject
to $CU=b$ is
\[
 U^*=R^{-1}C^T(CR^{-1}C^T)^+b,\qquad
 J^*=\tfrac12 b^T(CR^{-1}C^T)^+b,                                    \tag{3}
\]
where $^+$ is the Moore--Penrose inverse. If $b$ is outside that range,
repair is infeasible. Formula (3) is for signed controls; sign and budget
limits require the corresponding constrained quadratic program.
\end{theorem}
\begin{proof}
Set $y=R^{1/2}U$ and $D=CR^{-1/2}$. The problem minimizes
$\|y\|^2/2$ subject to $Dy=b$. In an orthogonal singular-value
decomposition $D=V\Sigma W^T$, feasibility sets each coordinate of
$W^Ty$ associated with a positive singular value equal to the corresponding
coordinate of $V^Tb$ divided by that value; zero singular values require
the corresponding coordinate of $V^Tb$ to vanish. The remaining free
coordinates uniquely minimize the squared norm at zero. Thus
$y^*=D^T(DD^T)^+b$, with squared norm $b^T(DD^T)^+b$.
Substitution yields (3) and the infeasibility statement.
\end{proof}

For $\delta\in(0,1)$, let $\eta$ be the declared marginal social value
of a unit of this log-capacity index. This is a linear welfare criterion,
not consumption utility inferred from the index. Since $A^h$ converges,
its powers are bounded and $\sum_h\delta^hA^h=(I-\delta A)^{-1}$.
Hence the discounted capacity gain from a finite control is exactly
\[
 \sum_{h\geq0}\delta^h\eta\rho_h
 =\eta\sum_{s=0}^{T-1}\delta^{s+1}
 w^T(I-\delta A)^{-1}Gu_s.                                         \tag{4}
\]
Indeed absolute convergence permits switching the finite control sum with
the infinite horizon sum, and the geometric series follows by telescoping
its finite partial sums; the remainder $\delta^{H+1}A^{H+1}$ tends to
zero. Subtracting the actual $U^TRU/2$ gives net welfare in this declared
criterion. Least-cost permanent repair need not maximize (4), because
transitory effects and the timing of repair also have value.

\section{Finite-horizon identification and its limits}
Fix baseline covariates $X$, contraction $D\in\{0,1\}$, an intermediate
history $H$ measured before policy, and policy $A_0\in\{0,1\}$.
Here $A_0$ is a treatment indicator, not the transition matrix $A$.
Let $Y_h(d,a)=\log K_h(d,a)$ and $H(d)$ be well-defined cluster
potential outcomes. Assume consistency, positive conditional treatment
probabilities, integrable outcomes, and sequential exchangeability:
\[
 (H(d),Y_h(d,a))\perp D\mid X,\qquad
 Y_h(d,a)\perp A_0\mid X,D=d,H.
\]
The second condition is understood within the $D=d$ potential history;
it excludes residual policy selection after conditioning on that history.

\begin{proposition}[The required sequential adjustment]
Under these conditions the regime mean is
\[
 m_h(d,a)=\int\!\int
 \E[Y_h\mid X=x,D=d,H=h',A_0=a]\,
 dP(h'\mid x,D=d)\,dP_X(x).                                      \tag{5}
\]
Therefore $\tau_h=m_h(1,0)-m_h(0,0)$ and
$\rho_h=m_h(1,1)-m_h(1,0)$ are identified for observed horizons.
\end{proposition}
\begin{proof}
Condition first on $X$ and $H(d)$. Sequential exchangeability and
consistency identify the conditional outcome mean by the observed
conditional mean in (5). Baseline exchangeability and consistency identify
the law of $H(d)$ given $X$ by $P(H\mid X,D=d)$. Iterated expectation
then gives (5). Differences of marginal means do not require a joint
cross-regime distribution.
\end{proof}
Conditioning on the same realized post-contraction history in both
contraction regimes would change the target. Migration must be handled
through a fixed population or a stated exposure mapping, and measurement
revisions must be harmonized before interpreting $Y_h$ as capacity.
Component decompositions of $w^Tz_h$ are index accounting, not identified
causal mediation effects.

\begin{proposition}[Exact finite observational equivalence]
For every finite $H\geq1$, an unrestricted post-shock accumulation class
contains two deterministic economies with identical regime outcomes through
$H$, a temporary initial demand shock, and different permanent effects.
This remains true even if their subsequent state transitions are linear
and time invariant, when state dimension is unrestricted.
\end{proposition}
\begin{proof}
The permanent economy has scalar state $z_0=-1$, transition $z_{h+1}=z_h$
and log-capacity deviation $z_h$. In the recovering economy use $H+1$
states, all initially $-1$, transition
$(z_1,\ldots,z_{H+1})\mapsto(z_2,\ldots,z_{H+1},0)$, and observe the
first state. The initial shock acts only at date zero in both economies.
The observed deviations equal $-1$ at dates $0,\ldots,H$ in both,
but the recovering deviation is zero thereafter whereas the permanent
one remains $-1$. Give policy zero capacity effect in both and set
$K_h=\exp(Y_h)$, so all joint regime data coincide through $H$ and
capacity is positive. Both transitions are time invariant; one has a
unit mode and the other is nilpotent. Their permanent contraction effects
are respectively $-1$ and zero.
\end{proof}
This construction does not defeat identification in a known, fully
observed, low-dimensional model. It pinpoints the extra content required
for extrapolation: state coverage, dimension restrictions, invariance and
information about the permanent projection $P$. Estimating these with
credible demand and policy variation, validating financing effects, and
allowing irreversible nonlinear exit remain unresolved parts of the target.

\section{Completion Estimate}
60\%

This is a post hoc, heuristic estimate for the full catalogue target.
The linear state model yields exact permanent damage and repair conditions, least cost intervention formulas, and a finite horizon observational equivalence showing why permanence is not generally identified. Credible demand and stimulus variation, financing effects, nonlinear exit mechanisms, and an empirically identified welfare model are still missing.

\begin{thebibliography}{9}
\bibitem{yellen} J. L. Yellen (2016).
\emph{Macroeconomic Research After the Crisis}, 14 October.
Federal Reserve Board speech, subsection ``The Influence of Demand on
Aggregate Supply'', especially the paragraph on reversing supply effects.
Primary text consulted 9 October 2026.
\url{https://www.federalreserve.gov/newsevents/speech/yellen20161014a.htm}.
\end{thebibliography}

\end{document}
